% arXiv: force the pdflatex path. All content is text and tables; there are no
% .eps files, and this must stay within the first few lines to take effect.
\pdfoutput=1

\documentclass[11pt]{article}

\usepackage[utf8]{inputenc}
\usepackage[T1]{fontenc}
\usepackage[margin=1in]{geometry}
\usepackage{mathptmx}
\usepackage[scaled=0.92]{helvet}
\usepackage{courier}
\usepackage{amsmath,amssymb,amsthm,mathtools}
\usepackage[numbers,sort&compress]{natbib}
\usepackage{booktabs}
\usepackage{array}
\usepackage{tabularx}
\usepackage{longtable}
\usepackage{enumitem}
\usepackage{caption}
\usepackage{microtype}
\usepackage{xcolor}
\usepackage{tcolorbox}
\tcbuselibrary{breakable}
\usepackage[colorlinks=true,linkcolor=blue!60!black,citecolor=blue!60!black,urlcolor=blue!60!black,breaklinks]{hyperref}
\usepackage{url}
\usepackage{fancyhdr}

\pagestyle{fancy}
\fancyhf{}
\fancyhead[L]{\small Difficulty Is Not the Index}
\fancyhead[R]{\small Preprint}
\fancyfoot[C]{\thepage}
\renewcommand{\headrulewidth}{0.3pt}
\setlength{\headheight}{14pt}

\definecolor{navy}{HTML}{17324D}
\definecolor{bluec}{HTML}{2B6CB0}
\definecolor{greenc}{HTML}{3A7D44}
\definecolor{orangec}{HTML}{C05621}
\definecolor{lightblue}{HTML}{EAF2F8}
\definecolor{lightgreen}{HTML}{EDF7ED}
\definecolor{lightorange}{HTML}{FFF4E6}

\newcolumntype{Y}{>{\raggedright\arraybackslash}X}
\setlist{itemsep=2pt,parsep=0pt,topsep=4pt}
\renewcommand{\arraystretch}{1.15}

\newtheorem{proposition}{Proposition}
\newtheorem{corollary}{Corollary}
\theoremstyle{definition}
\newtheorem{definition}{Definition}

\newtcolorbox{claimbox}[1]{breakable,colback=lightgreen,colframe=greenc,boxrule=0.7pt,arc=2pt,left=7pt,right=7pt,top=6pt,bottom=6pt,title=\textbf{#1},fonttitle=\normalsize,before skip=8pt,after skip=8pt}
\newtcolorbox{cautionbox}[1]{breakable,colback=lightorange,colframe=orangec,boxrule=0.7pt,arc=2pt,left=7pt,right=7pt,top=6pt,bottom=6pt,title=\textbf{#1},fonttitle=\normalsize,before skip=8pt,after skip=8pt}

\newcommand{\Td}{\ensuremath{T_\delta}}
\newcommand{\pstar}{\ensuremath{p^{\star}}}
\newcommand{\dver}{\ensuremath{d_{ver}}}
\newcommand{\drisk}{\ensuremath{d_{risk}}}
\newcommand{\kap}{\ensuremath{\kappa}}
\newcommand{\DLO}{\ensuremath{\mathrm{DL}\Omega}}

\title{\LARGE Difficulty Is Not the Index\\[6pt]
\large Three Repairs to the Delegation Frontier,\\
and One Reporting Form}

\author{Chengshuai Yang\\
\small\texttt{platformaigpt@gmail.com}}

\date{August 24, 2026}

\begin{document}
\maketitle

\begin{abstract}
The Delegation Intelligence framework proposes that delegation be measured rather than asserted, and supplies the machinery: a task-difficulty scale T0--T6, a human cognitive-intervention scale H0--H5, a level scale DL0--\DLO{}, a separation of cognitive from governance intervention, an intervention-depth measure CID0--CID6, ten anti-inflation rules, and a central object---the \emph{delegation frontier} $F_A(h,p)=\max\{T: S_A(T,h)\ge p\}$, the hardest task band a system completes with probability at least $p$ under intervention budget $h$ \citep{yang2026delegation}. We take that framework as primary. This paper is its companion, supplying propositions where it supplies apparatus, and it argues that the frontier as defined cannot carry the weight placed on it, for four reasons that are repairable.

\textbf{The frontier is indexed on the wrong variable.} The framework defines an eight-coordinate difficulty vector including verifiability and cost of error, warns that the vector must remain available, and then defines the frontier over the aggregated band. Two consequences follow. The set $\{T : S_A(T,h)\ge p\}$ is not downward closed, so its maximum does not summarize it, and the framework's reading of $F$ as ``the hardest band the system reliably completes'' assumes a closure the aggregation destroys. Worse, success is never observed directly; what is observed is the declared verifier passing. The two differ by the verifier's false-pass rate, which is governed by the verifiability coordinate the band discarded. \textbf{The frontier's measurement error is largest exactly where the class is hardest to check, and it errs upward.}

\textbf{The frontier prices failure at zero.} $p$ is a threshold on success probability with no term for what the complement costs. We show $p$ is not a free parameter of the benchmark: a class fixes a minimum reliability $\pstar(\tau)=\rho/(1+\rho)$, where $\rho$ is the cost of a failure relative to the value of a success. Where residual harm is unbounded, $\pstar=1$ and the class is not delegable at any capability. The framework's rule on escalation then stops being a rule and becomes a corollary.

\textbf{The acceptor is not a coordinate.} The framework requires independent verification in three places and then reports a profile with no field for who accepted, so two profiles differing in whether the system graded its own work are indistinguishable. This worsens up the ladder rather than improving: the fraction of acceptance events whose criterion the system itself authored rises with the level by construction, approaching one at \DLO{}, where protection against self-redefinition of success is asked for as advice rather than as structure.

\textbf{Its load metric is defined by an unobservable.} HCIL's denominator is the total cognitive effort a completed task \emph{required}---a counterfactual. We replace it with $L=r(\tau)(\Td+c_i)$, every term of which is a timestamp, and observe that the intervention scale is invariant to authorization latency, which in deployment dominates the cost.

We give the repaired reporting form, an episode record that carries it, and falsification conditions. Where the framework is better than the task-property account it is compared against---intervention depth, the cognitive/governance split, the acceptance criterion inside the task tuple, and delegation compression---we say so and adopt it.
\end{abstract}

\tableofcontents
\bigskip

\section{Introduction}

\subsection{The framework this repairs}

The Delegation Intelligence framework makes a claim worth defending: that ``autonomous agent'' mixes at least three questions---how hard the task is, how much human cognition stays inside the loop, and how often the result survives checking---and that a delegation claim is meaningful only as a conjunction of all three \citep{yang2026delegation}. Its apparatus is correspondingly serious. A delegated task is a tuple $\tau=\langle G,X,C,R,V,H\rangle$ carrying its own acceptance criterion $V$. Difficulty is externally calibrated, so a task does not become easy by being solved. Interventions are typed, and cognitive assistance is distinguished from governance approval, so that a required signature does not count against the machine's intelligence. Success requires a criterion declared independently of the system's self-report. Ten anti-inflation rules make the levels falsifiable.

None of that is in dispute here, and most of it is better than what it is being compared against.

\subsection{Why it needs repair at all}

The difficulty is not in the apparatus but in the single object the apparatus is compressed into. The frontier
\begin{equation}
F_A(h,p) \;=\; \max\{\, T : S_A(T,h) \ge p \,\}
\label{eq:frontier}
\end{equation}
is the framework's answer to ``how autonomous is it,'' and it is read as \emph{``the hardest task band the system reliably completes under the allowed human cognitive-intervention budget.''} Three of the framework's own commitments are in tension with that object.

It commits to difficulty as a \textbf{vector} and then maximizes over a \textbf{band}. It commits to \textbf{verified} success and then treats $S$ as observed. It commits to \textbf{independent acceptance} and then reports a profile with no field for the acceptor. Each tension is local, and each is repairable without abandoning the scale.

\subsection{Contributions}

\begin{enumerate}
\item \textbf{Non-closure} (\S\ref{sec:closure}). $\{T:S_A(T,h)\ge p\}$ is not downward closed, because $T$ aggregates coordinates that move $S$ in opposite directions. The maximum is therefore not a summary of the set, and the frontier should be reported per class cell rather than as a scalar band.
\item \textbf{Upward bias} (\S\ref{sec:bias}). The observed frontier is computed from verifier-pass rate, not success. The two differ by the false-pass rate, which is governed by the coordinate the band discarded---so the frontier over-reports most on the classes with the weakest verification.
\item \textbf{The reliability threshold is set by the class} (\S\ref{sec:pstar}). $\pstar(\tau)=\rho/(1+\rho)$, and where residual harm is unbounded, $\pstar=1$. This makes the irreversibility floor quantitative, and derives the framework's escalation rule instead of asserting it.
\item \textbf{The acceptor coordinate, and $\sigma$} (\S\ref{sec:accept}). A delegation claim must name who accepted; and the self-authored acceptance fraction $\sigma$ rises with the level by construction, which puts the acceptance problem at its worst precisely at DL6 and \DLO{}.
\item \textbf{An observable load metric} (\S\ref{sec:load}). $L=r(\tau)(\Td+c_i)$ replaces HCIL, and \Td{}---a quantity the intervention scale is invariant to---is added.
\item \textbf{A merged reporting form and episode record} (\S\ref{sec:form}), extending the framework's own.
\end{enumerate}

\subsection{What this paper does not claim}

\textbf{It is not an argument that the levels are wrong.} DL0--\DLO{} is a reasonable ladder and the transition tests are the right kind of object. The repairs are to the measurement, not to the ordering.

\textbf{It is not a claim that difficulty is irrelevant.} Difficulty determines whether a task can be done at all and drives the exception rate. The claim in \S\ref{sec:collapse} is narrower: difficulty is the wrong thing to \emph{index the frontier on}, because it is not what determines whether the result can be checked or undone.

\textbf{Nothing here is measured.} The numbers in \S\ref{sec:pstar} are worked illustrations. \S\ref{sec:limits} says what would make any of it a finding.

\section{The two accounts, and where they meet}

Two documents were written about delegation on the same day by different routes: the framework this paper repairs, which builds measurement apparatus from the human-factors tradition \citep{sheridan1978,parasuraman2000,endsley1999,klein2004}, and a task-property account arguing that delegation is capped by properties of the work rather than of the worker \citep{yang2026taskproperty}. They were written independently. The overlap is large enough to be worth stating, because the disagreements are legible only against it.

\textbf{Where they agree.} Delegation is not a single autonomy label and must be conditioned on the task. Self-reported completion is not completion. Governance approval is not cognitive assistance, and conflating them makes a well-governed system look stupid. Interventions differ in kind, and a rate that averages a rescue with a clarifying question is misleading in the direction of comfort. Levels are established over families with intervals, not over demonstrations.

\textbf{Prior art both accounts share.} Autonomy indexed by the conditions it holds under is J3016's operational design domain \citep{sae2021j3016}, and the observation that the same system occupies different autonomy levels in different contexts is the adjustable-autonomy literature's \citep{scerri2002}. Morris et al.\ separate capability from autonomy explicitly \citep{morris2023levels}, which is the published taxonomy closest to both accounts---and the one \S\ref{sec:closure} disagrees with, since its autonomy column is read off the system while the same system occupies several of its rows at once. Reliance as calibrated by experience of failure is Lee and See's \citep{lee2004}; the exception-driven interaction the H1 budget names, including when to interrupt, is Horvitz's \citep{horvitz1999}; and the interaction guidelines that govern how such an interruption should be presented are Amershi et al.'s \citep{amershi2019}.

\textbf{Where the framework is better}, and adopted here without modification:

\begin{table}[h]
\centering\small
\begin{tabularx}{\textwidth}{@{}lYY@{}}
\toprule
& The framework has & The task-property account has \\
\midrule
Intervention depth & \textbf{CID0--CID6}, an ordinal scale of the highest task decision a human supplied & ``exception kind matters more than rate,'' listed as an unmodelled limitation \\
Governance & An explicit split, with the rule that governance does not lower the level & An authority ceiling distinct from an achieved grade---the same idea, less operational \\
The task object & $V$ \textbf{inside the tuple}: the acceptance criterion is part of what is delegated & A four-coordinate class that does not carry its criterion at all \\
Environment change & \textbf{Delegation compression}: solving a hard class repeatedly converts it into an easy one & Nothing \\
\bottomrule
\end{tabularx}
\caption{Four places the framework is ahead of the account it is compared against.}
\end{table}

The last of these turns out to be where the two accounts meet, and \S\ref{sec:compression} argues it is the mechanism the task-property account was missing.

\textbf{Where they disagree} is the subject of \S\S\ref{sec:collapse}--\ref{sec:load}. In one sentence: the framework indexes the frontier on how hard the work is, and the task-property account holds that the binding coordinates are how cheaply a mistake is found and how cheaply it is undone---which are in the framework's difficulty vector, and gone from its difficulty band.

\section{Repair 1 --- the T-band collapse}\label{sec:collapse}

\subsection{The vector is defined, and then discarded}

The framework defines difficulty as a vector before it is a band,
\begin{equation}
d(\tau)=\langle d_{dec},\, d_{hor},\, d_{unc},\, d_{tool},\, d_{nov},\, \dver,\, d_{coord},\, \drisk \rangle
\end{equation}
and says the vector must remain available \emph{``so that a system cannot hide a narrow weakness behind an average difficulty score.''} This is exactly right, and the frontier does not honour it: $F_A(h,p)$ is a maximum over $T$, and $T$ is the aggregate. The warning is stated in the difficulty section and not carried into the formalism that uses it.

Two of the eight coordinates are the ones at issue. \dver{} is verification difficulty. \drisk{} is cost of error. They are the two the task-property account says do the work, and the two that have nothing to do with how hard the thinking is.

\subsection{The maximum is not a summary of the set}\label{sec:closure}

\begin{proposition}[Non-closure]\label{prop:closure}
The set $\{T : S_A(T,h)\ge p\}$ is not downward closed. Therefore its maximum does not entail that the system sustains $p$ on every lower band, and $F_A(h,p)$ is not ``the hardest band the system reliably completes.''
\end{proposition}

\emph{Argument.} Downward closure would require $S_A(\cdot,h)$ to be monotone decreasing in $T$. It is not, because $T$ aggregates coordinates whose effects on $S$ have opposite signs at fixed $h$. Raising $d_{dec}$ or $d_{hor}$ lowers $S$ only where planning or horizon is the binding constraint. Raising \dver{} lowers $S$ under every $h$ and for every agent, because it attacks the conjunct that makes a success \emph{verified}. A single scalar cannot preserve an ordering its components disagree about.\hfill$\square$

The counterexample is ordinary rather than contrived. Take two classes for one unchanged agent:

\begin{table}[h]
\centering\small
\begin{tabularx}{\textwidth}{@{}lYlll@{}}
\toprule
Class & Band, by the framework's operational characters & \dver{} & \drisk{} & $S$ at H1 \\
\midrule
Refactor a module under a good test suite & \textbf{T3} --- nontrivial decomposition, tool choice, state tracking, error recovery, replanning & low & low & high \\
Compose and send one external message & \textbf{T0--T1} --- one operation, or a familiar short task with a known procedure & \textbf{maximal} & \textbf{maximal} & not estimable \\
\bottomrule
\end{tabularx}
\caption{One agent, two classes. The second is nearer the bottom of the difficulty scale and nearer the bottom of what anyone will delegate.}
\end{table}

If $S$ clears $p$ on T3 and does not on T1, then $\max$ returns T3 while a lower band fails, and the sentence ``the hardest band it reliably completes'' is false of the system it describes.

\textbf{The repair} is to report the frontier per class cell rather than as a band. Let $\kap=\langle \dver,\drisk\rangle$ and let $T_{do}$ be the remaining six coordinates aggregated as before. Then $F_A(h,p \mid \kap)$ is downward closed in $T_{do}$ within a cell---the property the scalar was reaching for---and comparisons across cells are reported as what they are: different questions.

\subsection{The measurement error is a function of the coordinate that was discarded}\label{sec:bias}

This is the more serious half, because it biases in a known direction.

\begin{proposition}[Upward bias]\label{prop:bias}
$S_A(T,h)$ is not observed. What is observed is $S^V_A(T,h)$, the probability that the declared verifier passes. $S^V \ge S$, with the gap equal to the verifier's false-pass rate $\phi$. Since $\phi$ is governed by \dver{}, and \dver{} is aggregated into $T$, the frontier's error is largest exactly on the classes whose verification is weakest---and it is always upward.
\end{proposition}

\emph{Argument.} The framework's Rule 5 requires that success be established by a criterion declared independently of the system's self-report, which is right and is not the issue. Any such criterion is itself an instrument with a false-pass rate: a test suite with coverage below one, a held-out set with leakage, an inspector reading a document. Passing the instrument is the event recorded. So the observed frontier is $\hat{F}=\max\{T:S^V\ge p\}$, and $\hat{F}\ge F$ pointwise. The size of the gap is the property of the class the band deleted.\hfill$\square$

\begin{cautionbox}{Three consequences}
\textbf{A high frontier on a weakly verifiable class is the expected artefact}, not a surprising result. The framework's Rule 1 forbids establishing a high claim on low-difficulty tasks; there is no corresponding rule forbidding establishing it on low-\emph{verifiability} tasks, and that is the direction the error runs.

\textbf{Some classes are not in the frontier's domain at all.} Where no external criterion exists, $S^V$ is undefined rather than small, so the class silently drops out of the maximum. A frontier computed over the classes that happen to have verifiers---on a corpus where the unverifiable ones were excluded by the methodology rather than by a finding---reports a number that is high because of what is missing from it. These are the classes people most want a delegation answer for.

\textbf{The verifier belongs in the report.} A frontier is a statement about a system \emph{and its evaluator}. Reporting $\hat{F}$ without $\phi$, or without a statement that $\phi$ is unknown, is reporting an instrument reading as a quantity.
\end{cautionbox}

\textbf{The repair} is three fields: the verifier's identity, its estimated false-pass rate or an explicit \texttt{unknown}, and the count of classes excluded from the frontier for want of a criterion. The third is the one that keeps under-reporting from looking like capability.

\section{Repair 2 --- the frontier prices failure at zero}\label{sec:price}

Equation~\eqref{eq:frontier} contains no term for what happens in the $1-p$. Every failure costs the same as every other, which is to say nothing. The framework is not unaware of this---its Rule 8 says a correct escalation can be preferable to a confident failure---but the frontier has nowhere to put the difference, so the observation is carried as a rule alongside the formalism instead of falling out of it.

\subsection{\texorpdfstring{$p$}{p} is not a free parameter}\label{sec:pstar}

Write, for a class $\tau$: $V$ for the value of a verified success, $C_{det}$ for the cost of detecting a failure, $C_{undo}$ for the cost of undoing it, and $C_{res}$ for residual harm that no expenditure undoes. Let
\begin{equation}
\rho(\tau) \;=\; \frac{C_{det}+C_{undo}+C_{res}}{V}
\end{equation}
be the class's \textbf{loss ratio}. Delegating is worth doing when expected value is non-negative, $S\,V-(1-S)\,\rho V \ge 0$, which gives:

\begin{proposition}[The class sets the threshold]\label{prop:pstar}
\begin{equation*}
\pstar(\tau)\;=\;\frac{\rho(\tau)}{1+\rho(\tau)}
\end{equation*}
is the minimum reliability at which class $\tau$ may be delegated at all. Reporting $F_A(h,p)$ with $p$ chosen by the evaluator rather than by the class is a category error, and any $p<\pstar(\tau)$ reports a frontier over tasks that should not have been delegated at that reliability.
\end{proposition}

\begin{table}[h]
\centering\small
\begin{tabularx}{\textwidth}{@{}llrY@{}}
\toprule
Class & $\rho$ & $\pstar$ & A benchmark reporting at $p=0.90$ \\
\midrule
Draft a document & 0.1 & 0.09 & far stricter than the class requires---under-reports \\
Refactor under tests & 1 & 0.50 & stricter than required \\
Change production configuration & 30 & 0.968 & \textbf{reports a frontier the class does not permit} \\
Send an external message & unbounded & \textbf{1} & \textbf{the class is not in the domain at any capability} \\
\bottomrule
\end{tabularx}
\caption{The reliability a class requires, against the reliability a benchmark conventionally quotes. Illustrative values (\S\ref{sec:limits}).}
\end{table}

The convention of quoting frontiers at $p=0.90$ is therefore simultaneously too strict for cheap-to-undo classes and too lenient for expensive ones, and the error is not small: for production configuration, the reported and the permissible thresholds differ by most of the remaining probability mass.

\begin{corollary}[The irreversibility floor, quantified]\label{cor:floor}
As $C_{res}\to\infty$, $\pstar\to 1$. No attainable reliability delegates such a class, so it has a floor no capability removes. This is the quantitative form of the observation that verification after the fact is not a remedy: adding nines buys delegation only where residual cost is finite, and reliability and reversibility are therefore not interchangeable \citep{perrow1984}.
\end{corollary}

\subsection{Rule 8 becomes a corollary}

An escalation is an outcome with $S=0$ and $C_{det}=C_{undo}=C_{res}\approx 0$: the system did not complete the task and also did not incur the loss. It therefore does not enter the loss term at all, and costs only human load.

\begin{corollary}\label{cor:escalation}
A correct escalation dominates a confident failure on every class with $\rho>0$, and its cost is a load term rather than a reliability term. Frequent escalation therefore lowers the frontier at H0/H1 without lowering expected value.
\end{corollary}

Both halves of Rule 8 are recovered---including the second, which reads in the original as a concession and here as an accounting identity. This is the test of a repair: rules that had to be asserted alongside the formalism should follow from it.

\subsection{The repaired frontier}\label{sec:repaired}

\begin{equation}
F^{\star}_A(h)\;=\;\max\Big\{\,T \;:\; S^V_A(T,h)\,V-\big(1-S^V_A(T,h)\big)\rho V - L(T,h)\;\ge\;E_{sup}(T)\,\Big\}
\end{equation}
where $L$ is the human load of \S\ref{sec:load} and $E_{sup}$ is the value of doing the class under supervision. It reduces to the original when $\rho$ is constant across classes and $L$ is ignored---the condition under which the original is a good approximation---and states plainly where it is not.

\section{Repair 3 --- the acceptor is not a coordinate}\label{sec:accept}

\subsection{The requirement is in the prose and not in the profile}

The framework requires independent acceptance three times: a result counts as success only under a criterion declared independently of the system's self-report; self-reported completion is insufficient; and a benchmark should not ask the same model that performed the work whether it succeeded. All correct. The reporting form is
\begin{equation*}
\text{DL profile}=\{F(\text{H0},p),\,F(\text{H1},p),\,F(\text{H2},p),\,\text{HCIL},\,\text{CID dist.},\,\text{success CI},\,\text{cost},\,\text{latency}\}
\end{equation*}
which has no field for who accepted.

\begin{proposition}\label{prop:acceptor}
Two systems whose results were accepted by different loci---one by a separated verifier, one by the system itself---produce identical profiles. The reporting form therefore cannot express the difference between a measured level and a self-certified one, and a rule stated in the prose but absent from the form is enforced by the diligence of whoever fills it in.
\end{proposition}

The underlying result is that delegation of execution scales without limit while delegation of acceptance only transfers: if the accepting party is the performing party, the criterion sits inside the write set of the thing being judged, any criterion can be satisfied by adjusting it, and no observation distinguishes a met criterion from a moved one. Capability makes this worse rather than better, since a more capable performer is better at finding the adjustment \citep{yang2026taskproperty}.

\textbf{The repair is a field.} Every delegation claim carries $\alpha$, the acceptance locus:

\begin{table}[h]
\centering\small
\begin{tabularx}{\textwidth}{@{}llY@{}}
\toprule
$\alpha$ & Acceptor & Admissible? \\
\midrule
$\alpha_0$ & The performing system & \textbf{No}---this is an assertion, not a level \\
$\alpha_1$ & A test declared before the work, run by the performer & Yes, if the performer cannot edit the test \\
$\alpha_2$ & A separate process with its own credential and enforced write set & Yes \\
$\alpha_3$ & A party that did not build the system & Yes, and required for published claims \\
\bottomrule
\end{tabularx}
\caption{The acceptance locus. The requirement is on locus, not on species.}
\end{table}

Acceptance may be a machine, and must not be the same machine. One base model in two processes with separate credentials and enforced write sets satisfies the requirement. One model in one context told to now act as the reviewer does not---the roles are then two functions of one process, and the arrangement certifies itself.

The engineering half of this already exists in the companion agent-design framework, which states the rule as \textbf{evaluator not in write set}, gives write-set depths D0--D4, and makes D0---the agent that may change nothing about itself---the one fit to hold judgement over another's promotion \citep{yang2026agentdesign}. Its manifest carries \texttt{verifier: external\_required} and \texttt{promoter\_required} as declared fields, enforced by file permissions, separate services and container boundaries rather than by policy. The acceptance coordinate is implementable in the chassis today. It is missing from the scale, not from the build.

\subsection{The self-authored acceptance fraction, and why the ladder makes it worse}

There is a subtler failure, and it is the one that matters most for the top of the scale.

The task tuple carries $V$, so the delegated task's acceptance criterion is human-supplied. But from DL3 upward the system \emph{constructs the task graph}: it invents subgoals, and each invented subgoal carries an invented done-condition. Those conditions are acceptance events, and the system authored them. The human's $V$ binds the root of the tree and nothing below it.

\begin{definition}[Self-authored acceptance fraction]
$\sigma$ is the share of acceptance events in an episode whose criterion the system itself wrote.
\end{definition}

\begin{proposition}\label{prop:sigma}
$\sigma$ rises with the DL level by construction. At DL0 it is zero, because the human specifies each operation and its acceptance. At DL3 it is $1-1/N$ for an episode with $N$ acceptance events, since only the root criterion is given. At DL6 the system derives goals and priorities, so criteria at the project layer are self-authored too. At \DLO{} the system selects its own missions, so $\sigma\to 1$ unless something structurally prevents it.
\end{proposition}

This inverts the framework's own reassurance. It treats independent verification as a rule applying uniformly up the ladder; in fact the ladder is a description of \emph{handing over criterion authorship}, and the amount of self-certification inside an episode is a monotone function of the level being claimed. \DLO{} asks for ``explicit protection against self-redefinition of success'' in its evidence desiderata---which names this exactly, and asks for it as advice. Advice is a property of the day, not of the system.

\begin{corollary}\label{cor:sigma}
A DL6 or \DLO{} claim is unmeasurable, rather than merely hard to measure, unless $\sigma$ is bounded below one by structure: a criterion register written before the work and outside the system's write set, or a separated acceptor holding the pen on subgoal criteria as well as on the root.
\end{corollary}

\textbf{The repair} is to report $\sigma$ alongside the level, and to record for each acceptance event who authored the criterion and when relative to the work. The framework already requires intervention \emph{timing} to be recorded, on the grounds that information supplied before the system meets the relevant uncertainty may be hidden scaffolding. The same argument applies with the sign reversed: a criterion authored \emph{after} the work is a criterion fitted to it.

\section{Repair 4 --- the load metric is defined by an unobservable}\label{sec:load}

\subsection{HCIL's denominator is a counterfactual}

\begin{equation*}
\text{HCIL}=\frac{\text{human cognitive effort required}}{\text{total cognitive effort required by the completed task}}
\end{equation*}

The denominator is the effort a complete solution \emph{required}---a quantity about a counterfactual solution nobody performed. It is not merely hard to instrument; no procedure would observe it even in principle, because the task was completed once, jointly, and the decomposition into what the machine needed and what the task needed is not present in the record.

The framework says so, and offers six proxies: intervention count, human cognitive minutes, human-generated next-step decisions, corrections, rescues, and the fraction of branches whose strategy the human chose. Each is observable. But a quantity defined by an unobservable and reported through six proxies with no stated combination rule is a proxy set, and the name on it suggests a measurement.

There is a second problem. HCIL is a \emph{ratio} normalized by task difficulty, so a fixed amount of human help scores better on a harder task. That reintroduces exactly the conflation \S\ref{sec:collapse} removes.

\subsection{The observable is load, not a ratio}

\begin{equation}
L\;=\;r(\tau)\times\big(\Td+c_i\big)
\end{equation}
where $r$ is exceptions per task on the class, \Td{} is the latency between an exception being raised and the human acting, and $c_i$ is the human's per-intervention cost. Every term is a subtraction of two timestamps or a figure the human can report about their own clock. Nothing about a counterfactual appears. Below H2-equivalent budgets, $r$ degenerates to the action count and the identity reduces to per-action supervision, which is the correct limit \citep{amdahl1967}.

The weighting the ratio was reaching for is better served by CID, which the framework already defines: report $L$ per CID band rather than dividing by an unknown, so that one deep rescue is not averaged against twenty approvals. This keeps the framework's best measurement idea and drops its weakest.

\subsection{The intervention scale is invariant to the term that dominates}\label{sec:tdelta}

H0--H5 measures \emph{whether} human cognition is needed and at what frequency. It says nothing about how long the human takes. The two systems below score identically on the intervention scale and are not comparably delegable:

\begin{table}[h]
\centering\small
\begin{tabularx}{\textwidth}{@{}Ylll@{}}
\toprule
& Intervention rate & \Td{} & $L$ per task \\
\midrule
System P & one exception per task, H1 & 4 hours & 4 h \\
System Q & one exception per task, H1 & 30 seconds & 30 s \\
\bottomrule
\end{tabularx}
\caption{The intervention scale is a rate; the cost is a rate times a latency.}
\end{table}

A framework whose purpose is to make delegation operational should carry the term that decides whether a delegated task finishes today.

\Td{} is also the quantity bounding the iteration rate of any governed improvement loop \citep{yang2026twoaxisv2}, which makes it the one measurement two otherwise separate frameworks both turn on. It is cheap---timestamp the raise and the response---and on the one deployed system examined for this corpus it had never been recorded, because the table logging authorizations has no timestamp column \citep{yang2026tdelta}. The repair is a column, and the finding is that a quantity six documents call decisive has gone unobserved for want of one.

\section{Where the two accounts meet: delegation compression}\label{sec:compression}

The framework contains an idea the task-property account lacks and needs. \textbf{Delegation compression}: a T4 problem solved repeatedly becomes a T2 procedure once the system has built the tools, templates, checks and knowledge; delegation changes the task environment and not only the agent. The framework adds the right reporting discipline---record both the original and the post-tool difficulty, because silently relabelling the task as always easy erases the achievement.

Placed against \S\ref{sec:collapse}, this is more than a caveat. The task-property account holds that the lever on delegation is the class rather than the agent: add tests and a class moves, add one-command rollback and another does, neither requiring the agent to improve at anything. Compression is the same operation performed \emph{by the agent on its own task environment}---and the coordinates it moves are exactly \kap{}. A test the agent writes raises verifiability. A migration it makes reversible lowers residual cost. A checklist it proceduralizes lowers both.

\begin{claimbox}{The two accounts meet at compression}
The durable way for a system to raise its own frontier is not to become more capable but to move the \kap{} coordinates of the classes it works in. \textbf{Capability improvements raise $S$ within a cell; compression moves the cell.}

Two consequences for measurement. \textbf{Compression must be reported against \kap{}, not only against $T$}---``T4 became T2'' is the visible shadow of ``\dver{} fell,'' and the second is the durable claim. And \textbf{compression is the one route by which an agent legitimately raises its own frontier}, which makes it the thing to instrument if the question is whether a system is becoming more delegable rather than merely more accurate.
\end{claimbox}

\section{One reporting form, merged}\label{sec:form}

A delegation claim, stated fully:
\begin{equation*}
\text{DL claim}=\big\langle\,\text{DL}n,\;\kap,\;T_{do},\;\alpha,\;\sigma,\;F^{\star},\;\pstar,\;r,\;\Td,\;\text{CID dist.},\;\phi,\;\rho,\;\text{cost},\;\text{excluded classes}\,\big\rangle
\end{equation*}

\begin{table}[h]
\centering\small
\begin{tabularx}{\textwidth}{@{}llY@{}}
\toprule
Field & From & Answers \\
\midrule
DL$n$ & the framework & The compact label \\
$\kap=\langle \dver,\drisk\rangle$ & \S\ref{sec:collapse} & Which cell the claim is in \\
$T_{do}$ & the framework, less \kap{} & How hard the doing is \\
$\alpha$ & \S\ref{sec:accept} & Who accepted \\
$\sigma$ & \S\ref{sec:accept} & What share of criteria the system wrote \\
$F^{\star}$ & \S\ref{sec:repaired} & The frontier, net of failure cost and load \\
$\pstar$ & \S\ref{sec:pstar} & The reliability the class requires---not the one chosen \\
$r$, \Td{}, CID & \S\ref{sec:load} & What it costs the human, and how deep the help went \\
$\phi$ & \S\ref{sec:bias} & How much the verifier lets through, or \texttt{unknown} \\
$\rho$ & \S\ref{sec:pstar} & What a failure costs relative to a success \\
excluded classes & \S\ref{sec:bias} & What the frontier is silent about \\
\bottomrule
\end{tabularx}
\caption{The merged reporting form. Four fields and a denominator separate it from the original.}
\end{table}

The episode record, extending the framework's own with the fields the repairs require. Added blocks are marked:

\small
\begin{verbatim}
{
  "task_id": "tsk_...",
  "task_band": "T3",
  "difficulty_vector": {"horizon": 3, "uncertainty": 2, "verification": 3},
  "kappa": {"verifiability": 3, "reversibility": 1},          // added
  "intervention_budget": "H1",
  "human_interventions": [
    {"type": "approval", "cognitive": false, "minutes": 0.2,
     "cid": 0,
     "raised_at":    "2026-08-24T13:02:11Z",                  // added
     "responded_at": "2026-08-24T13:04:47Z"}                  // added
  ],
  "acceptance": {                                             // added block
    "locus": "a2_separated_process",
    "verifier_id": "verify-worker-03",
    "false_pass_rate": null,
    "events": 7,
    "self_authored_criteria": 5,
    "sigma": 0.71,
    "root_criterion_declared_at": "2026-08-24T12:40:00Z"
  },
  "loss": {"value": 1.0, "c_det": 0.1, "c_undo": 0.2,         // added block
           "c_residual": 0.0, "rho": 0.3, "p_star": 0.231},
  "agent_replans": 4,
  "verification": "pass",
  "result": "success",
  "latency_seconds": 4200,
  "compute_cost": 12.4
}
\end{verbatim}
\normalsize

The additions are five. \kap{} carries the two coordinates the band aggregates away. The timestamp pair yields \Td{} by subtraction. The acceptance block carries $\alpha$, $\sigma$, the verifier's identity and its false-pass rate---here \texttt{null}, which is the honest value and the common one. The loss block carries the terms \pstar{} needs. Every one is a timestamp, a count, or a figure the person delegating already knows. None requires estimating a counterfactual, which was the objection to the metric it replaces.

The framework's ten anti-inflation rules survive, with one added and one demoted.

\begin{claimbox}{Rule 11 --- verifiability is a test condition}
No frontier may be reported without the verifier's identity, its false-pass rate or an explicit \texttt{unknown}, and the count of classes excluded for want of a criterion. Rule 1 forbids establishing a high claim on easy tasks; this forbids establishing it on uncheckable ones, which is the direction the measurement error actually runs.

\medskip
\textbf{Rule 8 is demoted to a corollary.} It follows from Corollary~\ref{cor:escalation} and no longer needs to be a rule.
\end{claimbox}

\section{Falsification}

\begin{table}[h]
\centering\small
\begin{tabularx}{\textwidth}{@{}lY@{}}
\toprule
Claim & Refuted by \\
\midrule
Prop.~\ref{prop:closure}, non-closure & A system whose measured $S$ is monotone in $T$ across classes spanning a wide range of \dver{}---which would show the band preserves the ordering after all \\
Prop.~\ref{prop:bias}, upward bias & Verifier false-pass rates uncorrelated with \dver{} across a task corpus, or frontiers that do not shift when a stronger verifier replaces a weaker one on the same class \\
Prop.~\ref{prop:pstar}, \pstar{} & Delegation decisions tracking a fixed $p$ rather than $\rho/(1+\rho)$ once $V$, $C_{det}$, $C_{undo}$ and $C_{res}$ are elicited per class \\
Cor.~\ref{cor:floor}, the floor & An unbounded-$C_{res}$ class sustainably delegated at H0/H1 with outcomes matching reversible classes \\
Prop.~\ref{prop:acceptor}, the acceptor & Self-accepted work replicating under independent check at rates comparable to independently accepted work, on matched classes \\
Prop.~\ref{prop:sigma}, $\sigma$ & Episodes at DL3 and above whose self-authored acceptance fraction does not rise with level once acceptance events are individuated \\
\S\ref{sec:tdelta}, \Td{} & Human load failing to track $r\times(\Td+c_i)$ once all three are logged \\
\S\ref{sec:compression}, compression & Frontier gains that persist without any movement in \kap{}---capability alone moving the cell rather than the position within it \\
\bottomrule
\end{tabularx}
\caption{What would refute each claim.}
\end{table}

\textbf{Proposition~\ref{prop:bias} is the one most worth attacking}, because it says published frontiers are biased and names the direction. The experiment is cheap: take one class, measure the frontier under two verifiers of known differing coverage, and see whether the reported frontier moves with the instrument.

\section{Limitations}\label{sec:limits}

\textbf{Nothing here is measured.} The tables in \S\ref{sec:closure} and \S\ref{sec:pstar} are worked illustrations assembled from the pattern rather than from logged episodes---which by this paper's own standard makes them illustrations and not findings. Turning them into findings requires the episode record of \S\ref{sec:form} populated across classes, which is cheap and not done.

\textbf{$\rho$ is as hard to elicit as the thing it repairs.} $C_{res}$ in particular is a judgement about harm that does not undo, and people are poor at it. The claim is only that $\rho$ is elicitable in principle while HCIL's denominator is not, and that a badly estimated $\rho$ is visible in the report whereas a chosen $p$ hides the same judgement inside a convention.

\textbf{$\sigma$ requires acceptance events to be individuated}, and what counts as one event is a modelling choice. A system decomposing finely reports a higher $\sigma$ than one decomposing coarsely for the same behaviour. The ordering across levels is robust to this; the value is not.

\textbf{$\phi$ is rarely known.} Most verifiers have no false-pass estimate, which is why \S\ref{sec:form} admits \texttt{unknown} as a value. A field that is usually null is a weak instrument---its purpose is to make the absence visible rather than to supply the number.

\textbf{The repairs make the report heavier}, and there is a real tension between completeness and adoption. A form nobody fills in measures nothing. The minimum preserving the results is three fields: \kap{}, $\alpha$, and the two timestamps.

\textbf{\kap{} is used as a pair of scalars and is not.} Verification cost, latency and coverage are distinct, as are rollback cost, rollback latency and blast radius. A quantified version of \S\ref{sec:collapse} needs six coordinates where this paper uses two.

\textbf{The framework repaired here is an unpublished working paper}, as is the task-property account of \S\ref{sec:accept}. Every result is conditional on their premises, and the conditionality transfers.

\section{Conclusion}

The Delegation Intelligence framework is right that delegation should be measured as a conjunction of difficulty, intervention and reliability, and right that the alternative is a rhetorical label. The repairs here are to the object that conjunction is compressed into.

A frontier indexed on difficulty maximizes over an aggregate whose components disagree, so its maximum does not summarize the set it is drawn from; and because success is observed through a verifier, the frontier's error grows with the verifiability coordinate the aggregate discarded, always upward. A frontier that thresholds on success probability prices failure at zero, when in fact the class fixes the threshold---$\pstar=\rho/(1+\rho)$, unattainable where harm does not undo. A profile with no field for the acceptor cannot distinguish a measured level from a self-certified one, and the share of acceptance criteria the system itself wrote rises with the level, so the problem is worst at the top of the ladder rather than solved there. And a load metric whose denominator is a counterfactual should be replaced by one whose every term is a timestamp---including the latency the intervention scale is invariant to, and which nobody has recorded.

None of this requires abandoning T0--T6, H0--H5, DL0--\DLO{}, or the frontier. It requires the frontier to be indexed by the class rather than by the difficulty, priced against the cost of being wrong, stamped with who accepted, and costed in latency. Four fields and a denominator.

The practical form is a single sentence added to the framework's own. It asks: \emph{how difficult a goal can I hand to this intelligence, how little must I think for it, and how often does the result survive independent verification?} To which: \textbf{on which class, checked by whom, and at what cost when it is wrong.}

\bibliographystyle{plainnat}
\bibliography{references}

\end{document}
